%! Symmetric Positive Definite Matrices — Unit 6, Lesson 6.3 — Warm-Up (Answer Key)
\documentclass[10pt]{article}
\usepackage{linalg-article}
\usepackage{linalg-key}   % pulls in -boxes; do NOT also load -boxes
\newcommand{\TallMath}[1]{$\displaystyle #1\rule[-1.4em]{0pt}{3.2em}$}

\begin{document}

\pageheader{Unit 6, Lesson 6.3}{Warm-Up --- Answer Key}
\namedateperiod

\vspace{0.12in}

\begin{notesbox}{Getting Ready: Three Checks We'll Use Today}
\small Today's matrices are \emph{symmetric}, and their eigenvectors turn out perpendicular. These three
quick checks are the whole toolkit. Answer quickly.

\vspace{4pt}
\textbf{1. Transpose \& symmetry.} The \emph{transpose} $S^{\mathsf T}$ flips $S$ across its diagonal (rows
become columns). A matrix is \textbf{symmetric} when $S=S^{\mathsf T}$. Which of these is symmetric?
\[
  A=\begin{bmatrix} 2 & 1 \\ 1 & 2 \end{bmatrix}, \qquad
  B=\begin{bmatrix} 2 & 3 \\ 1 & 2 \end{bmatrix}.
  \qquad\text{Symmetric: }\ans{$A$ (its mirror entries match: $1=1$)}
\]

\vspace{4pt}
\textbf{2. Perpendicular?} Two vectors are perpendicular exactly when their dot product is $0$ (Unit~4).
Is $\begin{bsmallmatrix}1\\1\end{bsmallmatrix}$ perpendicular to $\begin{bsmallmatrix}1\\-1\end{bsmallmatrix}$?
\[
  \begin{bmatrix}1\\1\end{bmatrix}\cdot\begin{bmatrix}1\\-1\end{bmatrix}=\ans{$1-1=0$}
  \qquad\text{Perpendicular? }\ans{yes}
\]

\vspace{4pt}
\textbf{3. Make it a unit vector.} To \emph{normalize} a vector, divide by its length
$\|\mathbf{v}\|=\sqrt{\mathbf{v}\cdot\mathbf{v}}$. Find the length of $\mathbf{v}=\begin{bsmallmatrix}1\\1\end{bsmallmatrix}$
and write the unit vector in its direction:
\[
  \|\mathbf{v}\|=\ans{$\sqrt{2}$}, \qquad \frac{\mathbf{v}}{\|\mathbf{v}\|}=\ans{$\frac{1}{\sqrt2}\begin{bsmallmatrix}1\\1\end{bsmallmatrix}$}.
\]
\end{notesbox}

\begin{teachernote}
These three checks \emph{are} today's method, in order: \textbf{(1)} spotting a symmetric matrix
($S=S^{\mathsf T}$ --- note $A$ is symmetric though \emph{not} diagonal); \textbf{(2)} the dot-product
perpendicularity test, which will confirm that a symmetric matrix's eigenvectors sit at right angles; and
\textbf{(3)} normalizing, the move that turns those perpendicular eigenvectors
$\begin{bsmallmatrix}1\\1\end{bsmallmatrix},\begin{bsmallmatrix}1\\-1\end{bsmallmatrix}$ into the orthonormal
columns of $Q=\frac{1}{\sqrt2}\begin{bsmallmatrix}1&1\\1&-1\end{bsmallmatrix}$. Debrief quickly, then name the
plan: ``perpendicular eigenvectors normalize into an orthonormal $Q$, giving $S=Q\Lambda Q^{\mathsf T}$.''
\end{teachernote}

\end{document}
